\section{Taxonomia Proposta}

Para classificar os frameworks do corpus de forma uniforme, propomos uma taxonomia de seis dimensões: especificação, contexto, papéis, execução, validação e portabilidade. Essas dimensões foram escolhidas por aparecerem repetidamente, ainda que com nomes diferentes, nos estudos analisados, e correspondem aos campos estruturados extraídos de cada estudo. A taxonomia é apresentada aqui como referencial; sua validação empírica contra o corpus, incluindo as dimensões fortes e as fracas, é reportada na seção de Resultados. A Tabela~\ref{tab:taxonomia} resume cada dimensão por uma pergunta orientadora e por indicadores observáveis.

\noindent\textbf{Especificação.} Descreve como o framework transforma intenções humanas em artefatos que orientam o trabalho do agente. No extremo mais fraco, a especificação é apenas um prompt; no mais forte, é um conjunto versionado de requisitos, critérios de aceitação, planos, tarefas, arquitetura e políticas. O desenvolvimento orientado por especificação torna essa dimensão explícita ao tratar a especificação como contrato versionável que dirige a geração \cite{DeepakBabuPiskala2026}, e AgileGen a operacionaliza por meio de critérios de aceitação executáveis \cite{SaiZhang2025}. Reversa desloca a origem da especificação: em vez de partir de um produto novo, recupera especificações operacionais a partir de sistemas existentes \cite{Macedo2026}.

\noindent\textbf{Contexto.} É o mecanismo pelo qual o framework decide o que o agente deve saber antes de agir, o que inclui documentos, código, histórico, decisões arquiteturais, regras locais e evidências coletadas. A engenharia de contexto é tratada como disciplina própria para agentes que operam sobre repositórios reais \cite{SeyedmoeinMohsenimofidi2025}, e AutoCodeRover demonstra que a recuperação estrutural do contexto no código melhora a qualidade das correções geradas \cite{YuntongZhang2024}. Em Reversa, o contexto é extraído do legado e rotulado por confiança \cite{Macedo2026}.

\noindent\textbf{Papéis.} Definem especialização, autoridade e expectativa de comportamento. Frameworks multiagente distribuem tarefas entre personas como gerente de produto, arquiteto, desenvolvedor, testador e revisor, o que reduz a ambiguidade do prompt e cria expectativas de saída. MetaGPT é o caso mais explícito, ao internalizar procedimentos operacionais padrão e atribuir artefatos específicos a cada papel \cite{SiruiHong2023}, enquanto ChatDev e CodePori mostram que a decomposição por papéis já era uma linha ativa de pesquisa \cite{ChenQian2023,ZeeshanRasheed2024}.

\noindent\textbf{Execução.} Indica se o framework apenas produz orientação ou se também aciona ferramentas, edita código, roda testes e interage com interfaces. Ferramentas agênticas modernas tornam a execução uma dimensão central: o agente não apenas recomenda, mas modifica o ambiente. AutoDev e AutoCodeRover operam diretamente sobre o repositório, editando código e executando testes em tarefas reais de melhoria de programas \cite{MicheleTufano2024,YuntongZhang2024}, e Prompt Sapper materializa a execução como uma plataforma de produção de cadeias de IA \cite{YuCheng2023}.

\noindent\textbf{Validação.} Compreende os mecanismos para verificar se o que foi produzido está correto, completo e alinhado ao contexto, o que pode incluir testes automatizados, revisão humana, artefatos de evidência, traços auditáveis e gates de prontidão. O asseguramento baseado em traços articula contratos, testes e governança \cite{Paduraru2026}; a colaboração verificável entre assistentes torna auditável cada ato do fluxo \cite{Merchant2025}; e Reversa explicita confiança e lacunas nos artefatos recuperados \cite{Macedo2026}. A validação é crítica porque agentes podem produzir resultados plausíveis sem aderência real ao sistema.

\noindent\textbf{Portabilidade.} Descreve se o framework depende de um fornecedor, IDE, modelo ou formato específico. No corpus, esta é a dimensão mais frágil: ela aparece quase sempre pelo lado negativo, como risco de aprisionamento e de dependência de plataforma, e o aprisionamento só é nomeado explicitamente em um estudo \cite{SeyedmoeinMohsenimofidi2025}. Quase nenhum framework trata portabilidade como propriedade de projeto positiva e deliberada, o que justifica lê-la como dimensão predominantemente diagnóstica, conforme detalhado na seção de Resultados.

\begin{table}[h!]
\centering
\caption{Dimensões da taxonomia proposta.}
\label{tab:taxonomia}
\small
\begin{tabularx}{\textwidth}{p{2.8cm}Y Y}
\toprule
\textbf{Dimensão} & \textbf{Pergunta orientadora} & \textbf{Indicadores observáveis} \\
\midrule
Especificação & Como a intenção vira contrato de trabalho? & Specs, PRDs, planos, stories, tarefas, critérios de aceitação \\
Contexto & Como o agente sabe o que é relevante? & Repositório, docs, hooks, regras, memória, evidências, lacunas \\
Papéis & Quem decide, quem implementa e quem revisa? & Personas, agentes, skills, responsabilidades, autoridade \\
Execução & O framework age no ambiente ou apenas orienta? & Edição de código, comandos, testes, navegador, terminal, IDE \\
Validação & Como erros são detectados antes de virarem entrega? & Testes, checklists, gates, artefatos, revisão humana, confiança \\
Portabilidade & O processo sobrevive fora de uma ferramenta? & Múltiplas integrações, formatos abertos, aprisionamento, instalação local \\
\bottomrule
\end{tabularx}
\end{table}
